\section{Dimension parts and genuine global irreducible components}

Let $A$ be a closed analytic subset of a complex manifold $M$ of dimension $n$.
The global irreducible components are the closures of the connected components
of the regular locus, as in \cite{demailly-cadg}, II.5.3. A set has pure
geometric dimension $k$ when every one of its regular points has dimension $k$.
This characterization uses actual local manifold dimensions; it assigns no
extraneous dimension certificate. Write $A^{(k)}=\overline{\operatorname{Reg}_k(A)}$.

\begin{lemma}[Pure geometric dimension of a global component]
\label{hir:global-component-pure-geometric-dimension}
\uses{ana:regular-component,ana:pure-dimension,rem:global-components}
\source{demailly-cadg}
\dcref{II.5.3}
\lean{Hironaka.HironakaDimensionGlobalComponents.isPureDim_closure_component,Hironaka.HironakaDimensionGlobalComponents.exists_unique_isPureDim_of_mem_globalComponents}
\leanok
Every global irreducible component of $A$ has a unique pure geometric dimension
$k\leq n$.
\end{lemma}
\begin{proof}
A global component $B$ is the closure of the connected component $C$ through
some regular point $c$ of $A$, say of dimension $k$. All points of $C$ are
regular points of $A$ of dimension $k$, and they are also regular points of
$B$ of dimension $k$. In particular $B$ has a regular point of dimension $k$.
The closure $B$ is an irreducible analytic set. Pure dimensionality of an
irreducible analytic set therefore implies that every regular point of $B$
has dimension $k$. The local dimension at $c$ is at most $n$. If $B$ also has
pure dimension $j$, its regular point $c$ gives $j=k$, proving uniqueness.
\end{proof}

\begin{lemma}[Closed unions of global components]
\label{hir:closed-global-component-subunion}
\uses{rem:global-components}
\source{demailly-cadg}
\dcref{II.5.3}
\lean{Hironaka.HironakaDimensionGlobalComponents.isClosed_sUnion_globalComponents_subfamily}
\leanok
The union of any subfamily of the global irreducible components of $A$ is
closed in $M$.
\end{lemma}
\begin{proof}
Every global component is closed, and their family is locally finite. Any
subfamily remains locally finite. A locally finite union of closed sets is
closed: at a point outside the union, choose a neighborhood meeting only
finitely many members, and remove those finitely many closed sets from the
neighborhood. The resulting neighborhood misses the entire union.
\end{proof}

\begin{theorem}[The irreducible-component description of dimension parts]
\label{hir:dimension-part-global-component-union}
\uses{hir:global-component-pure-geometric-dimension,hir:closed-global-component-subunion,hir:embedded-resolution-dimension-parts}
\source{demailly-cadg}
\dcref{II.5.3}
\lean{Hironaka.HironakaDimensionGlobalComponents.dimensionPart_eq_sUnion_globalComponents}
\leanok
For every $k\in\mathbb N$, the dimension part $A^{(k)}$ is exactly the union
of the global irreducible components of $A$ of pure geometric dimension $k$.
No compactness, empty-interior, or positive-dimension assumption is required.
\end{theorem}
\begin{proof}
Let $S$ be that union. If $c$ is a regular point of $A$ of dimension $k$,
the closure of its connected component in the regular locus is a global
component of pure dimension $k$ containing $c$. Hence
$\operatorname{Reg}_k(A)\subset S$. The preceding lemma makes $S$ closed,
so $A^{(k)}\subset S$.

Conversely, let $B$ be a global component of pure dimension $k$, and write
$B=\overline C$ for the regular-locus component through a regular point $c$
of $A$ of dimension $j$. That same point is regular of dimension $j$ on $B$,
so pure dimension $k$ gives $j=k$. Every point of $C$ is then regular of
dimension $k$ on $A$, and therefore
$B=\overline C\subset\overline{\operatorname{Reg}_k(A)}=A^{(k)}$.
Taking the union gives $S\subset A^{(k)}$ and the claimed equality.
\end{proof}

\begin{lemma}[Geometric purity of the dimension part]
\label{hir:dimension-part-geometric-purity}
\uses{hir:embedded-resolution-dimension-parts,lag:regular-point-basics}
\lean{Hironaka.HironakaDimensionGlobalComponents.isPureDim_dimensionPart}
\leanok
Every regular point of $A^{(k)}$ has local manifold dimension $k$.
\end{lemma}
\begin{proof}
The regular points of dimension $k$ on $A^{(k)}$ are dense in $A^{(k)}$.
At a regular point $y$ of dimension $j$ on that set, choose a neighborhood
where all points of $A^{(k)}$ are regular of dimension $j$. Density supplies
a regular point of dimension $k$ inside this neighborhood. Uniqueness of
local manifold dimension at that point gives $j=k$.
\end{proof}


\begin{lemma}[Component membership in a dimension part]
\label{hir:global-component-dimension-part-criterion}
\uses{hir:global-component-pure-geometric-dimension,hir:dimension-part-global-component-union,lag:regular-point-basics}
\lean{Hironaka.HironakaDimensionGlobalComponents.isPureDim_iff_subset_dimensionPart_of_mem_globalComponents}
\leanok
A global irreducible component $B$ of $A$ is contained in $A^{(k)}$ if and
only if it has pure geometric dimension $k$.
\end{lemma}
\begin{proof}
The reverse implication follows from the union description. For the forward
implication, write $B=\overline C$ where $C$ is the regular-locus component
through a regular point $c$ of $A$ of dimension $j$. Since $c\in B\subset
A^{(k)}$, every neighborhood of $c$ meets $\operatorname{Reg}_k(A)$. Take a
neighborhood where all points of $A$ are regular of dimension $j$. A point
of $\operatorname{Reg}_k(A)$ in it has both dimensions $k$ and $j$, so $j=k$.
The component-purity lemma gives pure dimension $k$ for $B$.
\end{proof}

\begin{theorem}[Pure local decompositions of a dimension part]
\label{hir:dimension-part-local-decomposition}
\uses{ana:dimension-strata,lag:local-statement,hir:embedded-resolution-dimension-parts}
\source{demailly-cadg}
\dcref{II.4.24, II.5.3}
\lean{Hironaka.HironakaDimensionGlobalComponents.localDecomposition_dimensionPart,Hironaka.HironakaDimensionGlobalComponents.exists_localDecomposition_dimensionPart}
\leanok
At every point of $A^{(k)}$ there is a finite local analytic decomposition
of $A^{(k)}$ whose pieces all have dimension $k$. More precisely, from any
local analytic decomposition $A\cap U=\bigcup_i A_i$ with piece dimensions
$d_i$, retain exactly the indices with $d_i=k$. They give such a decomposition
on the same $U$, with the same holomorphic equations for each retained piece,
dense regular loci and pairwise relatively thin intersections.
\end{theorem}
\begin{proof}
The local stratum-closure theorem identifies
$A^{(k)}\cap U=\bigcup_{d_i=k} A_i$. The selected index set is finite, so
choose a bijection from a finite consecutive index set onto it. Transport
the retained pieces and their equations along this bijection. Their union
is the displayed dimension part. The neighborhood is unchanged, hence
remains open, contains the chosen point and lies in its preferred chart.
Holomorphicity of every equation, the equation descriptions of each piece,
purity and density of its regular locus remain exactly the corresponding
properties of the original piece. The assigned dimension is $k$ because
the index is selected. Distinct selected indices remain distinct original
indices by injectivity of the bijection, so the relatively thin-intersection
property is inherited as well. Finally any point of $A^{(k)}$ lies in $A$;
existence of the original analytic local decomposition at that point supplies
the asserted decomposition of $A^{(k)}$.
\end{proof}
